% 第 16 章。数字：checks/check_ch16.py；引文：data/snapshots/merc_facts.csv（逐字核对见
% data/fetch/verify_web_quotes.py；孟第 27 页与劳第 100、102 页的几段按扫描页读，见快照元数据）；
% 数据：fgr2009_1720_shares（不随仓库分发，先运行 data/fetch/fgr2009_1720_shares.py）、finrev_facts；图：figsrc/fig_ch16.py。
\chapter{重商主义、约翰·劳与重农学派}\label{ch:mercantilism}

\begin{chapterguide}
\guidequestion{约翰·劳的纸币实验为什么先繁荣、后崩溃？重商主义者和重农学派各自怎样理解财富？}
\begin{guidegoals}
\item 概括重商主义的贸易顺差观（Mun 1664）及其逻辑漏洞；
\item 重建约翰·劳体系（1716—1720）的债务置换与纸币扩张机制，并与南海泡沫比较；
\item 用“坎蒂隆效应”说明新增货币的分配后果；
\item 读懂魁奈《经济表》（1758）作为最早的循环流量模型的结构。
\end{guidegoals}
\guideprereq{\chref{ch:finrevolution}\enspace\nameref{ch:finrevolution}；\chref{ch:accounts}\enspace\nameref{ch:accounts}}
\end{chapterguide}

1720 年 1 月 5 日，一个苏格兰人当上了法国的财政大臣。约翰·劳（John Law，1671—1729）这时已经管着法国唯一的发钞银行，也是一家大公司的首脑。按 Velde 的描述，这家公司在顶点时归国家从前的债主所有，征收全部的税，拥有或管理着大部分海外殖民地，还垄断了全部海外贸易。它随意发行不兑现的纸币，这种纸币是唯一的法偿货币，还债时债主不能拒收\cite{velde2003}。

公司的股价以里弗尔（livre）计，里弗尔是法国的记账单位。1718 年，一股在 250 里弗尔左右；1720 年 1 月，涨到了 10,000 里弗尔。银行发出的纸币，从 1719 年 4 月底的 3,800 万里弗尔，增加到 1720 年 6 月底的 23.59 亿里弗尔\cite{velde2003}。到 1720 年 12 月初劳离开法国时，按白银折算，股价已经跌回到 1718 年的水平附近（图~\ref{fig:merc_shares}(a)）。

同一年，伦敦的南海公司做了类似的事。它的股价从 1 月的约 120 英镑涨到 7 月的 950 英镑，10 月初跌到约 290 英镑\cite{garber1990}。五十多年后，亚当·斯密（Adam Smith，1723—1790）写道，可以把纸币增发到几乎任何数量的想法，“是所谓密西西比计划的真正基础；这是一个银行与股票投机兼而有之、也许是世界上从未见过的最荒唐的计划”\cite{smith1776}。

劳的计划背后有一套关于货币与财富的想法，回应的是一个更老的问题。一个国家的财富是什么，从哪里来？17 世纪的英格兰商人托马斯·孟（Thomas Mun，1571—1641）说，财富是贸易顺差带回来的金银。劳在 1705 年说，货币的多少决定了能雇多少人，而货币不必是金银。在巴黎做银行生意的理查德·坎蒂隆（Richard Cantillon，约 1680—1734）追问，新增的钱先到了谁手里。法国宫廷医生弗朗索瓦·魁奈（François Quesnay，1694—1774）则认为，只有土地在补偿了一切成本之后还留下一个剩余，他称之为“净产品”。他还把一年里的产品与货币怎样在三个阶级之间流转，画成了一张表。

劳的纸币为什么先繁荣、后崩溃？重商主义者与重农学派，又各自怎样理解财富？答案要从路易十四留下的债说起。

\section{当时的世界}\label{sec:merc-world}

劳与南海公司要处理的，都是国债。1715 年的法国与 1720 年的英格兰各欠了多少钱，欠的是什么样的债？金银在两国的政策里又占什么位置？

\subsubsection{路易十四留下的债务}
1715 年 9 月，路易十四（Louis XIV，1643—1715 年在位）去世。年幼的路易十五（Louis XV，1715—1774 年在位）即位，由奥尔良公爵（Philippe d'Orléans，1674—1723）摄政。Velde 重建了这个月的账目，国王的债有三类。永续年金（rentes）的债主永远不能要求偿还本金，这类债的本金合计 10.686 亿里弗尔。官职是向国王买来的官位，国王每年付给任职者一笔“俸给”，所以实际上也是一种债，合 8.30 亿里弗尔。流动债务是 9.189 亿里弗尔。三项合计 28.175 亿里弗尔，是当年收入 1.656 亿里弗尔的 17 倍（\chref{ch:finrevolution}表~\ref{tab:finrev-four}）。永续年金与官职的利息共 8,670 万里弗尔，占了收入的一半以上；扣除利息以前的基本盈余只有 4,770 万里弗尔\cite{velde2003}。

政府的对策是赖账。1715—1717 年，它对永续年金、官职的俸给与流动债务一再部分违约、降息，又把流动债务换成票面利率 4\% 的国家票据（billets d'État）。这种票据发出不久，就折价 37\%\cite{velde2003}。劳面对的就是这样一个局面，债务大，信用坏，利息吃掉了一半的收入。

\subsubsection{英格兰：不能赎回的年金}
英格兰作为债务人可靠得多，债务的形式却很乱。按 Garber 的数字，1720 年英国的国债约为 5,000 万英镑。其中 1,830 万在三家最大的公司手里，英格兰银行 340 万、东印度公司 320 万、南海公司 1,170 万。私人持有的可赎回国债有 1,650 万。约 1,500 万是\textbf{不可赎回的年金}，其中长期年金的期限在 72 到 87 年之间，短期年金是 22 年\cite{garber1990}。南海公司本身就是 1711 年把短期债务转换成股份的产物（\chref{ch:finrevolution}）。不可赎回的年金是一个难题。利率下降了，政府却不能像荷兰省那样按面值赎回、再以低息重借（式~\eqref{eq:finrev-perp}），只能等它们到期，或者说服债主自愿交换。1720 年，南海公司提出替政府做这件事。

\subsubsection{金银的位置}
两国的政策都把金银看得很重。按斯密的记述，西班牙与葡萄牙是欧洲主要金银矿的所有者，它们或者以最严厉的刑罚禁止金银出口，或者对出口课以可观的关税。苏格兰的一些旧法令，以重罚禁止把金银带出王国。法国与英格兰从前也有同样的政策\cite{smith1776}。为什么一个国家要把金银留在国内？

\section{孟与贸易差额}\label{sec:merc-mun}

这个问题最有名的答案，出自一个替出口白银辩护的人。重商主义者说的“财富”是什么？为什么贸易顺差被当成致富之道，这个想法又错在哪里？

\subsubsection{东印度公司的辩护人}
孟的书《英格兰得自对外贸易的财富》（\emph{England's Treasure by Forraign Trade}）是他的儿子在 1664 年出版的。按 1895 年重印本编者阿什利（William Ashley，1860—1927）的说法，书大约写于 1630 年\cite{mun1664}。阿什利在前言里介绍，孟在 1615 年当选为成立不久的东印度公司的董事。公司出口贵金属引起了抗议，为了替公司辩护，孟在 1621 年出版了《论英格兰与东印度的贸易》\cite{mun1664}。公司到东方去买胡椒与香料，要运去白银。在一个禁止或限制金银出口的世界里，这是一门需要辩护的生意。《英格兰得自对外贸易的财富》把这番辩护推广成了一套关于国家财富的理论。

孟的规则写在第二章。在他看来，增加国家的财富与金银，通常的办法是对外贸易，要遵守的规则是“每年卖给外国人的，在价值上要多于我们消费他们的”。他举了一个例子。假如英格兰每年出口价值 220 万英镑的货物，从外国进口 200 万英镑的货物供本国使用与消费。孟说，那样“我们可以确信，王国每年会富 20 万英镑，这笔钱必然以同样多的金银流入”\cite{mun1664}。写成式子，就是
\begin{equation*}
\text{金银的净流入}=\mathit{EX}-\mathit{IM},
\end{equation*}
即出口 $\mathit{EX}$ 减去进口 $\mathit{IM}$，孟的例子里是 $2.2-2.0=0.2$（百万英镑）。这里假定国与国之间除了买卖货物没有别的往来，没有运费与保险费的收支，也没有借贷，差额只能用金银结清。孟自己的账比这个式子宽。书的第二十章讲怎样编全国的贸易差额，出口要在原价上加上船的运费、航程的保险与商人的利润。进口则按在海外购买之处、连同全部费用装上本国船只时的成本计价，不按本国的估价\cite{mun1664}。今天的国际收支账（\chref{ch:external}）里，货物的差额只是经常账户的一部分，顺差也可以变成对外国人的债权，不一定变成金银。

\subsubsection{国家的账与商人的账}
孟的论证里最有意思的，是他区分了两本账。他说：“我们必须始终把王国的收益与商人的利润区分开来。”王国为一磅胡椒付出的，只是外国人从我们这里收到的价钱。商人除了这个价钱，还要付运费、保险费、关税和其他费用，在这样远的航程里，这些费用非常大。可是，“在王国的账上，这些都不过是我们自己人之间的交换，并没有减少王国的财富”\cite{mun1664}。

他用胡椒算了一笔账。假设胡椒在本国一直值每磅 2 先令，也就是 24 便士。先令与便士都是比英镑小的英国货币单位。从阿姆斯特丹的荷兰人那里买，商人每磅可以付到 20 便士，还能赚钱；到东印度去买，他最多付 3 便士\cite{mun1664}。孟说，这是一个巨大的好处，不只对本国自用的那一部分，也对每年从这里运往其他国家、以更高价格出售的大量胡椒\cite{mun1664}。

\begin{example}{王国的收益与商人的利润}\label{ex:merc-mun}
用孟的胡椒价格：本国每磅 24 便士，阿姆斯特丹 20 便士，东印度 3 便士\cite{mun1664}。教学假设：从阿姆斯特丹进货的费用可以不计；去东印度的远航，每磅胡椒要付运费、保险费与关税合计 18 便士，都付给英格兰的船主、保险人与国王。（1）在本国出售时，商人走哪条路赚得多？（2）每磅在本国消费的胡椒，两条路分别要付给外国人多少？（3）如果商人把东印度买来的胡椒按阿姆斯特丹的价格卖给外国人，王国每磅得到多少金银？
\solution （1）阿姆斯特丹：$24-20=4$ 便士；东印度：$24-3-18=3$ 便士。商人宁愿去阿姆斯特丹。（2）走阿姆斯特丹，王国付给荷兰人 20 便士；走东印度，只付给印度人 3 便士，其余 18 便士付给了本国人。按孟的算法，王国每磅省下 $20-3=17$ 便士的金银。商人的利润少了 1 便士，王国的收益多了 17 便士，这就是孟要区分两本账的原因。（3）转卖时，王国每磅付出 3 便士、收回 20 便士，净流入 17 便士的金银，与运费由谁收无关。孟的算法有一个隐含的假定：那 18 便士所雇的船、水手与资金，如果不用在远航上，什么也不会生产。如果它们本来可以做别的事，远航的代价就不只是付给外国人的 3 便士。
\end{example}

\subsubsection{播种与收获}
既然国家的账只看与外国人的收支，孟就可以为出口白银辩护。在他看来，出口货币与出口货物一样，只要是为了贸易，就会增加王国的金银\cite{mun1664}。他打了一个比方：“如果我们只看农夫在播种时的举动，看他把大量的好谷子撒进地里，我们会把他当成疯子，而不是农夫；但是想到他在收获时的劳动，那才是他辛劳的目的，我们就会看到他这番举动的价值与丰厚的收成。”\cite{mun1664}斯密后来概括说，孟把对外贸易的这种做法比作农业的播种与收获\cite{smith1776}。

斯密把商人们反对禁令的理由归纳为两条。第一，为购买外国货而出口金银，不一定减少国内的金银，因为这些货物可以再出口到外国高价卖出，带回比当初运出去的多得多的金银。第二，禁令阻止不了金银出口，因为金银体积小而价值大，很容易走私出去。斯密说，这些论证“一部分是可靠的，一部分是诡辩”，可它们说服了听的人，也就是议会与君主的咨议会、贵族与乡绅\cite{smith1776}。

结果，在法国与英格兰，禁止出口的金银只限于本国的铸币，外国铸币与金银块的出口是自由的。在荷兰与另一些地方，连本国铸币也可以自由出口。政府的注意力从守住每一块金银，转到了盯住贸易差额。斯密的评语是，它“从一种徒劳的关心，转到另一种复杂得多、麻烦得多、同样徒劳的关心上”。孟的书名成了政治经济学的一条基本准则，不只在英格兰，在所有商业国家都是如此\cite{smith1776}。

\begin{compare}{金银禁运与贸易差额}{禁止金银出口}{孟的贸易差额}
\cmprow{谁主张}{西班牙、葡萄牙；苏格兰的旧法令；从前的法国与英格兰}{孟（约 1630 年写成，1664 年出版）与东印度公司的商人}
\cmprow{管什么}{每一笔金银的出入}{一年里全部对外贸易的差额}
\cmprow{出口金银}{禁止，或课以重税}{为了贸易可以出口，像播种一样}
\cmprow{能不能做到}{金银体积小、价值大，容易走私}{只能靠关税与贸易管制间接影响}
\cmprow{共同的前提}{一国的财富以它拥有的金银来衡量}{同左}
\cmpsum{孟把注意力从每一块金银移到了整个贸易差额上；他驳倒了禁运，却保留了“金银就是财富”的前提。}
\end{compare}

\begin{pitfall}[出口金银一定让国家变穷]
这是孟要驳倒的说法，至今仍常听到。按孟自己的账，出口白银到东印度买胡椒，再把胡椒卖给外国人，收回的白银多于付出的白银；只看第一笔交易，就像只看农夫撒种。今天的说法换成了“对某国的贸易逆差就是损失”。一国对一个国家有逆差，可以对另一个国家有顺差；对外支付的钱买回了货物，货物也是财富。判断一国的得失，要看全部的交易，也要看换回来的是什么，而不是看某一笔钱流向了哪里。
\end{pitfall}

\subsubsection{漏洞在哪里}
孟的论证有四个漏洞，后来的人一一指了出来。第一，金银不等于财富。一个国家年年顺差、积攒金银，等于年年把货物送给外国人，只换回金属。金属本身不能吃也不能穿，它的用处在于能买东西，斯密对重商主义的整个批评就从这里出发（\chref{ch:smith}）。第二，国内的资源不是免费的。例~\ref{ex:merc-mun} 的 18 便士，在国家的账上是“自己人之间的交换”，可它雇用的船与人本来可以生产别的东西。第三，不是所有国家都能同时顺差。世界上所有国家的贸易差额加起来等于零，一国的顺差就是别国的逆差；人人都照孟的规则做，规则就不成立了。

第四，顺差不能持久。金银流进一个国家，那里的物价会上涨，出口变贵、进口变便宜，顺差自己就会缩小。孟自己也注意到，王国里钱多了，本国的货物就会变贵；货物贵了，就用得少，卖得少。可他由此得出的结论，是赚来的钱要继续用在贸易里，不要囤起来，不是顺差会自己消失\cite{mun1664}。坎蒂隆讨论货币的那一章（\secref{sec:merc-cantillon}），也描述了矿山的新钱抬高物价、外国货变得便宜、金银流往外国的过程\cite{cantillon1755}。把这些观察表述成一个自动调节国际收支的机制的，是休谟（David Hume，1711—1776），叫作价格—铸币流动机制（\chref{ch:bullion}）。孟的顺差和休谟的机制里，流动的都是金银。

\section{劳的货币理论}\label{sec:merc-law}

在劳看来，货币不必是金银。他在 1705 年为什么主张多发纸币，说这能让国家变富？他心目中的纸币以什么为担保？

\subsubsection{缺钱的苏格兰}
劳在 1705 年于爱丁堡出版了《论货币与贸易》（\emph{Money and Trade Consider'd}），向苏格兰议会建议设立一家发行纸币的银行。斯密后来写道，“苏格兰的产业因为缺少货币来推动而萎靡不振，这是著名的劳先生的看法”。劳提议设立一种特别的银行，斯密说，他似乎以为这家银行可以发行与全国土地的全部价值相等的纸币。苏格兰议会没有采纳。后来，当时的法国摄政奥尔良公爵有所变通地采纳了它\cite{smith1776}。

劳的出发点是：“国内贸易取决于货币。货币多，雇的人就比货币少时多。”好的法律可以让货币流通得更充分，可是在他看来，没有一条法律能让它走得更远；没有更多的钱来流通、支付更多人的工资，就不能让更多的人工作。人们也可以靠信用工作，只要信用能流通，“那么信用就是货币，对国内与对外贸易的作用都一样”。所以，增加货币，不论雇主赚不赚钱，都增加了国民财富\cite{law1705}。

这把孟倒了过来。孟说，顺差带来金银；劳说，“更多的货币雇用更多的人，就会有剩余可以出口”\cite{law1705}。货币在先，顺差在后。

\subsubsection{以土地为担保的纸币}
劳并不认为货币本身有价值。他说：“货币不是用来交换货物的那个价值，而是货物借以交换的那个价值。”\cite{law1705}那么用什么做货币最好？劳认为是土地。在他看来，“土地生产一切，白银只是产品之一”。白银或任何别的产品的数量都可能增减，土地却不增不减。土地还可以用纸来转让，所以比白银更具备货币所需的其他性质\cite{law1705}。

他设想的计价方法是，一英亩土地每年付 100 单位，按 20 年的租金计算，值 2,000 单位\cite{law1705}。这正是\chref{ch:finrevolution}的永续年金公式，式~\eqref{eq:finrev-perp}。每年付 100、价值 2,000，对应 $100/2000=5\%$ 的利率。他的结论是：“所提议的纸币，数量始终与需求相等，于是人民有活干，国土得到改良，制造业进步，国内外贸易得以开展，财富与力量也随之而来。”\cite{law1705}

用\chref{ch:silver}的数量方程，式~\eqref{eq:silver-qty}，可以看清劳的主张。$MV=PY$ 里，货币数量论假定 $Y$ 不受 $M$ 的影响，多出来的钱只会抬高物价。劳假定的是另一种情形。有大量闲置的人与土地，$M$ 增加时 $Y$ 跟着增加，$P$ 不必上涨。按 Garber 的转述，劳的货币理论是在存在未利用资源的环境里提出的。在这种环境里，发行纸币会永久地扩大实际的商业，对新货币的需求随之增加，足以让物价不受压力\cite{garber1990}。

按 Velde 的说法，劳一直相信，金属货币低劣而浪费，最好由纸币，或者由一种流动性很高、带利息的证券来代替\cite{velde2003}。Garber 还引了熊彼特（Joseph Schumpeter，1883—1950）的评价：“他把自己计划的经济学推敲得如此出色，甚至可以说深刻，足以让他跻身有史以来第一流的货币理论家之列。”\cite{garber1990}

\begin{intuition}[劳把孟倒了过来]
孟与劳都认为货币的多少关系到国家的财富，分歧在因果的方向与货币的形式。孟说，卖多买少，金银才会流进来；劳说，先有足够的货币，人才有活干，才有剩余可以出口。孟的货币只能是金银，要从外国人那里挣来；劳的货币可以是纸，国家自己就能发行。这个想法在有失业的时候有几分道理，这是后来凯恩斯（John Maynard Keynes，1883—1946）问题的一个早期版本。它的危险在“与需求相等”这几个字。需求是在什么价格、什么利率下的需求？如果银行按固定的条件向所有想借钱的人放款，对纸币的需求可以无限地增长。斯密后来讨论的“真实票据”原则，即银行只给因货物真实卖出而开出的票据放款，也面对同一个问题（\chref{ch:smith}）。
\end{intuition}

\section{劳的体系（1716—1720）}\label{sec:merc-system}

Velde 说，劳的公共财政实验从 1716 年 5 月设立通用银行（Banque générale）开始，到 1720 年 12 月他逃离法国为止\cite{velde2003}。这几年里，劳怎样一步步把法国的国债换成公司的股份，又把金银换成纸币？为什么钉住股价会让纸币失控？

\subsubsection{银行}
银行一开始是私人的。它发行 1,200 股，每股 5,000 里弗尔，大部分可以按面值用国家票据支付，而国家票据当时的市价比面值低 60\%。买新发的股份叫认购，认购的人不必一次付清，起初只需付价款的四分之一。银行最初的资产是 375,000 里弗尔现金与 112.5 万里弗尔国家票据，所以最初买一股只要 690 里弗尔现金\cite{velde2003}。算一下，$1200\times5000/4=150$ 万里弗尔，正好是两项资产之和。每股先付的 1,250 里弗尔里，现金 312.5 里弗尔，国家票据面值 937.5 里弗尔。国家票据按面值的四成计，约值 375 里弗尔，合计约 690 里弗尔。

1717 年 4 月 10 日的法令规定，个人可以用银行的纸币纳税\cite{velde2003}。按 Velde 的说法，银行平均每年发行 4,000 万到 5,000 万里弗尔的纸币，同时保持着合理的铸币准备，约为 50\%。1718 年 12 月 4 日，摄政以国王的名义按面值买下全部股份，银行改为王家银行（Banque royale），这时流通中的纸币是 3,950 万里弗尔。这件事“第一次给了国王一台能用的印钞机”。此后，银行的纸币很快从一家私人银行的负债、代表固定数量白银的债权，变成了唯一的法偿货币，与任何标准都脱了钩\cite{velde2003}。

\subsubsection{公司}
1717 年 8 月，劳获准设立西方公司（Compagnie d'Occident），开发路易斯安那\cite{velde2003}。这块殖民地就是密西西比河的流域。到 1717 年，殖民地的人口只有约 500 人，指的是殖民者，不算当地的印第安人\cite{velde2003}。Garber 说，用“密西西比”称呼劳的体系，就是从这门生意来的\cite{garber1990}。公司首次发行 200,000 股，每股 500 里弗尔，全部按面值用国家票据支付。认购在 1718 年 12 月 31 日结束，共卖出 1 亿里弗尔，其中 40\% 归国王所有，是用印好而没有花掉的国家票据买的\cite{velde2003}。Velde 指出，对国家票据的持有人来说，认购就是把一张 4\% 的债券换成一家公司的股份，而公司的主要资产就是同样的债券与路易斯安那\cite{velde2003}。

此后公司不断并购。1718 年 8 月 1 日，它以每年 402 万里弗尔取得烟草专卖。1719 年 5 月，它买下东印度公司与中国公司，整个集团改组为印度公司（Compagnie des Indes），按 Garber 的说法，新公司垄断了法国在欧洲以外的全部贸易\cite{garber1990}。7 月 25 日，公司以 5,000 万里弗尔取得九年的铸币权。8 月 27 日，它取得了总包税所（Fermes Générales）。包税，就是由包税人承包征税，交给国家一个固定的数。总包税所征收法国大部分的消费税，约占政府收入的 30\%，租约每年 5,200 万里弗尔，此前是 4,850 万。8 月下旬，公司又买断了征收全部直接税的官职，直接税约占收入的 55\%\cite{velde2003}。

为了付这些钱，公司继续发行股份。1719 年 6 月，它发行 50,000 股，每股 550 里弗尔。7 月又发行 50,000 股，每股 1,000 里弗尔。9—10 月发行 300,000 股，每股 5,000 里弗尔\cite{velde2003}。认购 6 月的新股，要凭四股旧股，旧股于是被称为“母亲”，新股是“女儿”。7 月的新股要凭四个母亲与一个女儿才能认购，称为“孙女”。Velde 说，这个规定助长了旧股二级市场的狂热\cite{velde2003}。三次发行加上首发，共 60 万股。

表~\ref{tab:merc-system} 把体系的主要步骤排在一起。

\begin{table}[htbp]\centering\small
\caption{劳的体系，1716—1720 年}\label{tab:merc-system}
\begin{tabularx}{\linewidth}{@{}l Y@{}}
\toprule
日期 & 事件 \\\midrule
1716 年 5 月 & 通用银行成立，股份大部分用折价的国家票据支付 \\
1717 年 4 月 10 日 & 银行的纸币可以纳税 \\
1717 年 8 月 & 西方公司获准设立；首发 20 万股，每股 500 里弗尔 \\
1718 年 12 月 4 日 & 银行收归国王，改为王家银行；流通纸币 3,950 万里弗尔 \\
1719 年 5—8 月 & 并购东印度公司与中国公司（改称印度公司），取得铸币权、总包税所与直接税 \\
1719 年 8 月 27 日起 & 公司提出借给国王 12 亿（后增至 16 亿）里弗尔，年息 3\%，用来偿还国债 \\
1719 年 9—10 月 & 发行 30 万股，每股 5,000 里弗尔；后来只收国债的偿还凭证，不收现金 \\
1719 年 12 月 29 日 & 宣布 1720 年每股股息 200 里弗尔，股价 9,413 里弗尔 \\
1720 年 1 月 5 日 & 劳任财政大臣 \\
1720 年 2 月 27 日 & 私人持有的金银铸币不得超过 500 里弗尔，100 里弗尔以上的支付必须用纸币 \\
1720 年 2 月 & 王家银行并入公司 \\
1720 年 3 月 5 日 & 公司以每股 9,000 里弗尔的固定价格买卖股份；金属里弗尔贬值三分之一 \\
1720 年 5 月 21 日 & 法令：股价分期降到 5,000 里弗尔，纸币降到面值的一半；5 月 27 日撤销 \\
1720 年 7 月 17 日 & 银行停止兑付 \\
1720 年 12 月初 & 劳离开法国；此后纸币与股份都交到 Visa 清算，纸币大多换成了国债 \\
\bottomrule
\end{tabularx}
\datasource{Velde（2003），第 10—38 页。}
\end{table}

\subsubsection{债务置换}
1719 年 8 月 27 日起，公司向国王提出一笔永久贷款，年息 3\%，用来偿还国王的债务。贷款起初是 12 亿里弗尔，9 月 17 日与 10 月 10 日增加到 16 亿。每年 3,600 万（后来增加到 4,800 万）里弗尔的利息，实际上从公司每年应付总包税所的 5,200 万租金里扣除\cite{velde2003}。国王的债大都可以这样偿还。按 Velde 的说法，永续年金与官职按其法律性质包含一个看涨期权，也就是说，债主永远不能要求偿还本金，债务人却可以随时偿还\cite{velde2003}。1719 年 8 月 31 日下令偿还的债务，本金合计 14.998 亿里弗尔，其中市政厅的永续年金 9.617 亿。到 1720 年 7 月，两批合计 19.771 亿\cite{velde2003}。

公司拿什么借给国王？卖股份。8 月 26 日，偿还国债的消息公布以前，股价为 3,600 里弗尔；9 月 9 日，涨到 5,350 里弗尔。两天后，公司请求政府允许它按每股 5,000 里弗尔卖股份，筹集 5 亿。发行很成功，9 月 26 日与 10 月 2 日又按同一价格发行，合计 10 亿\cite{velde2003}。$5+10=15$ 亿，正是 30 万股乘以 5,000 里弗尔。后来，股份不再收现金，只收国库发给债主的偿还凭证，以及政府的其他无记名债务；公司那笔 3\% 的债券始终没有发行\cite{velde2003}。所以 Velde 说，“这一操作不过是一次巨大的置换：把平均利率 4.5\% 的国债换成公司的股份。公司的利润来自政府欠它的 3\% 的利息，加上商业与包税的利润”\cite{velde2003}。

图~\ref{fig:merc_system} 画出了前后的变化。体系以前，税由各种包税人与收税官征收，他们交给国家一个固定的数。国家每年要付给债主约 9,000 万里弗尔，剩下的才能用于政府的支出。体系之下，公司包揽了全部的征税，又插在国家与债主之间。它付给股东不少于 4,800 万里弗尔的可变股息，国家有了更多的钱可以花\cite{velde2003}。约 9,000 万里弗尔的利息，与 19.77 亿里弗尔的本金相符，$19.77\times4.5\%\approx0.89$ 亿。

\fig{merc_system}{劳的体系以前与之下的法国财政。看两半张图里国家付出的利息：上半部是体系以前，包税人与收税官征税 $\tau$，交给国家一个固定的数 $\bar\tau$，国家每年付给债主约 9,000 万里弗尔，政府支出 $G$ 不能超过 $\bar\tau$ 减去 9,000 万；下半部是体系之下，公司征收全部的税，交给国家 $\bar\tau$ 减去 4,800 万（国家欠公司的 3\% 利息），付给股东不少于 4,800 万的可变股息，还经营海外贸易，政府支出的上限变成 $\bar\tau$ 减去 4,800 万。示意图，按 Velde（2003）图 6 重画。}{fig:merc_system}

国王什么时候合算？设国王原来欠债 $D$，平均利率 $i$。公司发行 $N$ 股、每股价格 $P$，收进国王的债券，$NP=D$；股东得到的是公司的收益 $E$，其中有国王付的 3\% 的利息，也有包税、铸币与贸易的利润，这些本来大都归国王或包税人。对国王来说，这等于以每年 $E$ 的代价借到了 $D$，成本是 $E/D$。置换比原来的债便宜，当且仅当
\begin{keyeq}
\frac{E}{NP}<i\qquad\Longleftrightarrow\qquad\frac{P}{E/N}>\frac{1}{i},\label{eq:merc-swap}
\end{keyeq}
其中 $E/N$ 是每股收益，$P/(E/N)$ 是市盈率。第二个式子是把第一个式子两边取倒数，不等号随之反向。按 $i=4.5\%$，市盈率要高于 $1/0.045\approx22$。Velde 也说，要给国王比他原来的债更好的条件，公司股票的市盈率必须高于 22\cite{velde2003}。

\begin{intuition}[股份是一种利息可变的国债]
对国王来说，股份像一种不必付固定利息的债：公司赚得多，股东拿得多；赚得少，损失由股东承担，国王不必违约。股价越高，换掉同样多的国债需要的股份越少，这种“债”的成本就越低。劳要股价高，不只是为了让股东高兴，也是为了让国王借得便宜。Garber 说，在劳的计划下，“法国实际上把国库私有化了”\cite{garber1990}。
\end{intuition}

\begin{example}{置换合不合算}\label{ex:merc-swap}
Velde 的表 6 记下了公司宣布的股息与当时的股价：1719 年 7 月 26 日，宣布 1720 年上半年每股股息 60 里弗尔，当时 25 万股，股价 1,960 里弗尔；1719 年 12 月 29 日，宣布 1720 年全年每股 200 里弗尔，当时 60 万股，股价 9,413 里弗尔\cite{velde2003}。（1）把宣布的股息当作每股收益，两次的价格股息比各是多少？按式~\eqref{eq:merc-swap}，置换对国王合算吗？（2）12 月的股价隐含的“借款成本”是多少？（3）按 12 月的股息，60 万股每年要付多少？劳 1723 年算过，付 200 里弗尔的股息给 40 万股需要 8,000 万里弗尔的收入\cite{velde2003}，这与 60 万股差在哪里？
\solution （1）7 月：$1960/60\approx33$；12 月：$9413/200\approx47$；与表 6 一致。两者都高于 22，按市场价格，置换对国王合算。（2）$200/9413\approx2.1\%$，不到原来 4.5\% 的一半；劳自己说他的目标利率是 2\%\cite{velde2003}。（3）$200\times600{,}000=1.2$ 亿里弗尔，比体系以前每年付给债主的约 9,000 万还多。这与第（1）问并不矛盾：式~\eqref{eq:merc-swap} 假定股份的市值正好等于换掉的国债，$NP=D$；可是 60 万股按 9,413 里弗尔值约 56 亿里弗尔，是要偿还的约 20 亿国债的近三倍。股份大多是按 5,000 里弗尔或更低的价格发出去的；按 5,000 里弗尔算，200 里弗尔的股息对应的价格股息比是 25，只比 22 高一点。劳的 8,000 万（$200\times400{,}000$）没有算公司作为贷款抵押品持有的 10 万股，以及国王持有的同样多的股份\cite{velde2003}。股价高时置换便宜，前提是公司真能挣出它宣布的股息；这要等到\secref{sec:merc-test}再检验。
\end{example}

\subsubsection{纸币与钉住的股价}
股份在涨，纸币也在涨。1718 年 12 月，通用银行的纸币为 4,000 万里弗尔。按 Velde 的表 4，王家银行的纸币 1719 年 8 月底为 3.806 亿里弗尔，1720 年 2 月底为 10.697 亿\cite{velde2003}。纸币是怎样发出去的？按 Velde 的说法，银行看来只是凭公司司库签字的借据，把纸币借给公司；1720 年 2 月，银行干脆并入了公司\cite{velde2003}。当时人估计，法国的金银铸币存量约为 12 亿里弗尔。1719 年 12 月 1 日的法令说，核准发行的 6.4 亿纸币足以满足“流通与一切商业活动”；三个月之内，这个数就增加到了三倍\cite{velde2003}。

同时，金银被一步步赶出流通。1720 年 2 月 27 日起，任何人持有的金银铸币不得超过 500 里弗尔，100 里弗尔以上的支付只能用纸币。3 月 11 日，又宣布全面废除金银货币。4 月 1 日，合同里的金银条款一律作废\cite{velde2003}。图~\ref{fig:merc_notes} 是 Velde 整理的纸币数量。

股价也有人托着。早在 1719 年 10 月 5 日，公司就命令司库以 5,000 里弗尔收购市场上的任何股份。这等于保证每个债主都能按面值拿回纸币，Velde 称之为“债务货币化的第一个迹象”\cite{velde2003}。1719 年 12 月下旬到 1720 年 2 月中，公司的收购处买进了 8 亿里弗尔的股份，约为市值的 16\%，货币也增加了同样多。当时有 30 万股原始股，公司按 9,925 里弗尔收购。另有至多 30 万份认购权证，每份值 6,600 里弗尔。两者的市值合计 49.6 亿里弗尔\cite{velde2003}。$8/0.16=50$ 亿，两个数相符。2 月 22 日，公司从国王手里买回他的 10 万股，每股 9,000 里弗尔；到 3 月 1 日，市价从 9,925 里弗尔的托市价跌到 8,500 里弗尔\cite{velde2003}。

3 月 5 日，公司开设了另一个办事处，按每股 9,000 里弗尔的固定价格买卖股份\cite{velde2003}。这一天，劳还把金属里弗尔贬值了三分之一，新铸的银币面值从 1 里弗尔改为 1.5 里弗尔，同时承诺“银行券是一种不会有任何变动的货币”\cite{velde2003}。银行券就是银行发出的纸币。钉住股价以后，谁想把股份换成纸币，都可以按 9,000 里弗尔换。设固定价格为 $\bar p$，公众卖给公司的股份为 $\Delta S$，银行为此新发的纸币就是
\begin{keyeq}
\Delta M=\bar p\,\Delta S .\label{eq:merc-peg}
\end{keyeq}
纸币的数量不再由银行决定，而由持股人决定。他们想要多少现钱，银行就得印多少。Garber 说，这一政策“实际上把股份变成了面额 9,000 里弗尔的银行券”\cite{garber1990}。从 3 月到 5 月下旬，公司又花了 12.135 亿里弗尔，买进了 27\% 的股份\cite{velde2003}。按式~\eqref{eq:merc-peg}，$12.135\text{ 亿}/9000\approx13.5$ 万股。同一时期，表 4 的纸币从 3 月 5 日的 10.899 亿增加到 5 月 31 日的 22.351 亿，增加了 11.452 亿，与公司收购股份花的钱大致相当\cite{velde2003}。

\begin{intuition}[钉住股价就是放弃控制货币]
钉住汇率的央行必须按固定价格买卖外汇，于是放弃了对货币数量的控制（\chref{ch:mundell}）。劳钉住的是自己公司的股价，道理一样。股份随时能按 9,000 里弗尔换成纸币，持有股份就等于持有纸币，另外还有股息；人们一旦觉得股份不值 9,000，就会卖给公司，纸币就会成倍增加。股价钉住了，代价是纸币的价值不再钉得住。
\end{intuition}

\fig{merc_notes}{劳的纸币，1719—1721 年。看实线与褐色虚线什么时候分开：实线是发出减去销毁的纸币（Velde 表 4 第 (1) 列），从 1719 年 4 月底的 3,800 万里弗尔增加到 1720 年 6 月底的 23.59 亿里弗尔；褐色虚线是估计的流通量按 1719 年的白银里弗尔折算的价值（第 (4) 列），1720 年 5 月底最高，为 16.255 亿，1721 年 1 月 8 日只剩 4,660 万。折算率是每马克白银 60 里弗尔，这里的马克是称金银的重量单位，不是德国的货币。两条线在 1720 年 3 月分开，是因为 3 月 5 日金属里弗尔贬值了三分之一；6 月以后越分越开，是因为纸币对铸币折价，加上一部分纸币已经收回但没有销毁。浅色带是钉住股价的 3 月 5 日至 5 月 21 日。数据：Velde（2003）表 4。}{fig:merc_notes}

\begin{case}{巴黎 · 1720 · 5 月 21 日的法令}
3 月 5 日的承诺只维持了两个半月。5 月 21 日，劳宣布分期降低股份与纸币的价值：股价从 9,000 里弗尔逐月降到 5,000 里弗尔，纸币到 12 月 1 日降到面值的一半\cite{velde2003}。按 Garber 的说法，劳认为自己把股价定得太高了；这一步完成以后，到 12 月，“只会剩下 23 亿里弗尔的纸”，其中纸币 13 亿、股份 10 亿\cite{garber1990}。股价没有等到逐月下调。按 Frehen、Goetzmann 与 Rouwenhorst 整理的数据，印度公司的股价 5 月 16 日还是 9,005 里弗尔，28 日是 6,075 里弗尔，30 日只剩 4,000 里弗尔\cite{fgr2009}。5 月 27 日，5 月 21 日的法令被撤销\cite{velde2003}；Garber 说劳本人在 5 月底被免职，很快又复职，主持了此后的紧缩\cite{garber1990}。

承诺一旦违背，就收不回来了。按表 4 的估计，流通中的纸币按 1719 年的白银计，5 月 31 日值 16.255 亿里弗尔，6 月 30 日值 11.903 亿，7 月 31 日只值 5.152 亿\cite{velde2003}。7 月 17 日，银行停止兑付\cite{velde2003}。12 月初，劳离开了法国。他至少有 950 万里弗尔的财产投在法国的房地产上，Velde 说，对一个打算尽快脱身的人来说，“这可不是什么好主意”\cite{velde2003}。劳似乎真的相信自己的体系。
\end{case}

\subsubsection{清算}
体系的善后叫作 Visa，即对纸币与公司证券的登记与核减，印度公司的股份也交到了 Visa\cite{velde2003}。按 Velde 的说法，发出的 28 亿里弗尔纸币中，22 亿最后进了 Visa 的清算，变成了国债。交到 Visa 的 22.112 亿里弗尔被核减到 17 亿，平均核减 23.5\%，但各人之间差别很大。500 里弗尔以下的债权一律不减，这些小额债权占了人数的一半、金额的 40\%；其余 60\% 的金额平均核减 39\%\cite{velde2003}。$60\%\times39\%\approx23.4\%$，与平均数相符。

对约 200 个在体系里赚了“过多利润”的人，另征了一笔税，债务又减了 1.9 亿。剩下的约 15 亿换成了 4\% 的终身年金与 2\% 的永续年金\cite{velde2003}。Visa 发出的清算凭证，1722 年在市场上约值面值的 25\%。Velde 估计，一张纸币的普通持有人最后拿回了面值的约 20\%，这也是 1720 年 11 月初纸币的市价\cite{velde2003}。同一年的伦敦，一家不发钞的公司也在用股份收购国债。

\section{南海公司}\label{sec:merc-southsea}

在 Garber 看来，南海公司比劳的公司简单。它没有大规模收购商业公司，也不承担铸币、征税或发行纸币这样的政府职能。英国与南美的贸易实际上被西班牙人挡住了，所以对这段经济史来说，要紧的只是它持有的国债\cite{garber1990}。它怎样用股份收购英格兰的国债？它的股价为什么涨得这样快，又跌得这样快？

\subsubsection{收购国债的权利}
斯密记下了南海公司做过的生意。第一门是凭《乌得勒支和约》给它的“阿西恩托”（Assiento）专营权，向西属美洲贩运非洲奴隶。人们不指望这门生意能赚多少钱，作为补偿，公司又获准每年派一艘许可船直接去西属美洲贸易。斯密说，这艘许可船“十次航行，据说只有 1731 年王家卡罗琳号那一次赚了不少，其余几乎每次都多少亏了钱”\cite{smith1776}。

1720 年，南海公司取得了收购非公司持有的 3,100 万英镑国债的权利，为此同意付给政府至多 750 万英镑。收购的条件按债的种类定。公司每收进每年 100 英镑的长期年金，可以增发面值 2,000 英镑的股份；收进同样的短期年金，可以增发 1,400 英镑。可赎回国债每 100 英镑面值，增发 100 英镑。政府为公司收进的国债付息，1727 年以前年息 5\%，此后 4\%\cite{garber1990}。为了让法案通过，公司向议会要员与国王的宠臣行贿，合计 130 万英镑。Garber 说，这些贿赂“给这段插曲添了几分阴暗，但它们本身并不是即将发生欺诈的信号”\cite{garber1990}。转换以前，南海公司的股份面值为 1,170 万英镑，到投机结束时增加到 2,280 万英镑\cite{garber1990}。

关键在于转换的比率没有写死。按 Garber 的说法，法案只规定了公司可以增发多少股份，股份与国债的交换比率由公司自由决定。公司给交换的股份定了远高于面值的价格，于是剩下一些核准而未用的股份，可以自由地卖出去\cite{garber1990}。股价越高，公司用越少的股份就能换到同样多的年金，剩下的配额越多。

\begin{example}{南海公司的报价}\label{ex:merc-offer}
1720 年 5 月 19 日，公司给每年 100 英镑长期年金的持有人开出的条件是：面值 700 英镑的股份（7 股），加上 575 英镑的债券与现金。报价时南海公司股价约 400 英镑；转换以前，这样一笔年金的市价约为 1,600 英镑\cite{garber1990}。（1）报价值多少？（2）股价跌到多少，持有人接受报价才会吃亏？（3）公司为这笔年金可以增发面值 2,000 英镑的股份，给出 700 英镑以后，剩下的配额按 400 英镑的股价值多少？
\solution （1）$7\times400+575=3{,}375$ 英镑，是年金市价的两倍多，与 Garber 的数字相同。（2）$7P+575=1600$，$P=(1600-575)/7\approx146$ 英镑。Garber 说，只要股价不跌到 146 英镑以下，年金持有人就不会吃亏，“所以这个报价极有吸引力”\cite{garber1990}。（3）还剩面值 1,300 英镑，即 13 股，按 400 英镑值 5,200 英镑。股价越高，第 (1) 问的报价越诱人，第 (3) 问留给公司的也越多：公司和年金持有人都盼着股价上涨。
\end{example}

\subsubsection{从 950 英镑到 290 英镑}
股价确实在涨。1720 年 1 月，每 100 英镑面值的股份约值 120 英镑；3 月 21 日，转换法案在议会首次获得通过，股价从约 200 英镑跳到 300 英镑\cite{garber1990}。6 月 17 日，公司第三次发行现金认购的股份，面值 500 万英镑，即 5 万股，每股市价 1,000 英镑。股价随即从 745 英镑跳到 950 英镑\cite{garber1990}。6 月 24 日到 8 月 22 日，公司为了派发股息关闭了过户登记，这段时间的成交实际上是远期交易，即现在成交、以后交割。Neal 认为顶点是 7 月 1 日的 950 英镑，Scott 记为 1,050 英镑，Garber 说后者大概包括了宣布的 10\% 的股票股息\cite{garber1990}。通过所有的转换，南海公司收进了公众持有的不可赎回年金的 80\%、可赎回国债的 85\%\cite{garber1990}。

6 月，议会通过了《泡沫法》（Bubble Act），禁止未经特许设立公司，也禁止已有的公司把特许状扩展到未经许可的新业务上。8 月 18 日，法令对南海公司的几个竞争者执行，这些公司的股价立刻受到压力\cite{garber1990}。南海公司自己的股价从 8 月 31 日的约 775 英镑，跌到 10 月 1 日的约 290 英镑。登记在册与将要发给公众的股份共 212,012 股，8 月 31 日的总市值约 1.64 亿英镑。一个月里，市值“蒸发了约 1.03 亿英镑，超过原来那笔沉重国债价值的两倍”\cite{garber1990}。$212{,}012\times775\approx1.64$ 亿；$212{,}012\times(775-290)\approx1.03$ 亿。善后时，议会没收了公司董事与几位官员的财产，共 200 万英镑，交给公司，又免除了公司为取得转换权而答应付给政府的 710 万英镑\cite{garber1990}。

这一年的巴黎与伦敦还有一个共同点，市面上的钱多了。多出来的钱先到了谁手里，又让谁吃了亏？

\begin{case}{伦敦 · 1720 · 英格兰银行的挤兑}
南海公司崩溃时，英格兰银行也遭到挤兑。坎蒂隆在《论一般商业的性质》里记下了银行的办法\cite{cantillon1755}。银行先是多派几个职员给持票人点钱，大笔的钱用六苏与十二苏的小银币来点，以拖延时间。苏是比里弗尔小的法国货币单位。一部分钱付给在那里等了一整天的普通持票人，“但最大的几笔付给朋友，他们把钱拿走，再偷偷送回银行，第二天重演一遍”。这些办法还不够时，银行开出认购书，邀请有信誉、有偿付能力的人联合起来，为大笔款项作保，维持银行券的信用与流通。坎蒂隆说，1720 年英格兰银行的信用就是靠这最后一着保住的：“公众一知道认购书已经由有钱有势的人签满，就不再挤到银行去，存款也照常送来了。”\cite{cantillon1755}止住挤兑的，不是银行金库里的金银，而是人们相信有足够多的人会站在银行背后。这还不是后来的最后贷款人。最后贷款人是中央银行在恐慌中向别的银行放款（\chref{ch:goldstandard}），这里却是一群有钱人联合起来，为一家银行作保。两者靠的都是让持票人相信，银行背后的钱足够多。
\end{case}

\begin{compare}{劳的体系与南海公司}{劳的体系（法国）}{南海公司（英格兰）}
\cmprow{换掉的国债}{约 20 亿里弗尔；Velde 说公司用股份替换了“全部国债”}{约 3,100 万英镑非公司持有的国债里的大部分（不可赎回年金的 80\%、可赎回国债的 85\%）；全部国债约 5,000 万}
\cmprow{公司还做什么}{包税、铸币、海外贸易；银行并入公司，发行纸币}{只持有国债；对西属美洲的贸易几乎无利可图}
\cmprow{货币}{钉住股价，纸币随之增发；金银被废除}{不发钞；英格兰银行与金银铸币照旧}
\cmprow{股价}{约 250 里弗尔涨到 10,000 里弗尔（1720 年 1 月），按白银计又跌回约 300}{约 120 英镑涨到 950 英镑（7 月），又跌到约 290 英镑（10 月 1 日）}
\cmprow{善后}{Visa 清算，纸币持有人平均拿回面值的约 20\%；劳出走}{没收董事财产 200 万英镑，免除公司欠政府的 710 万英镑}
\cmpsum{两者都是用股份换国债；劳同时把股份变成了钱，所以他的崩溃也是货币的崩溃；南海公司的崩溃主要是股价的崩溃，虽然也引发了英格兰银行的挤兑。}
\end{compare}

\fig{merc_shares}{两家公司的股价，1718—1720 年。(a) 看实线与褐色虚线之间的距离：实线是印度公司（1719 年 5 月以前为西方公司）的股价，以当时的里弗尔计（1720 年起主要是纸币），1718 年 8 月为 273.75 里弗尔，1720 年 1 月 8 日最高，为 10,100 里弗尔，3 月 5 日至 5 月 21 日钉在 9,000 里弗尔附近；褐色虚线从 3 月 5 日开始，按 Velde 表 4 各个日期（多为月末）的折算率把股价折成 1719 年的白银里弗尔：3 月 5 日的 9,000 里弗尔到 3 月底只值 6,000，11 月 25 日的最后一个报价只值约 300，回到了 1718 年的水平。浅色带是钉住股价的时期。(b) 看两条线的落差：实线是南海公司的股价，褐色虚线是英国东印度公司的股价，都是每 100 英镑面值的价格；南海公司从 1719 年 8 月的 113 英镑涨到 1720 年 7 月的 950 英镑，12 月初跌到 185 英镑，东印度公司最高只到 420 英镑。日期都按格里历（伦敦当时用的儒略历晚 11 天）。数据：Frehen、Goetzmann 与 Rouwenhorst（2009）的 1720 年股价数据库（印度公司一栏据 Velde 与 Murphy），Velde（2003）表 4。}{fig:merc_shares}

\section{坎蒂隆：新钱先到谁手里}\label{sec:merc-cantillon}

坎蒂隆说，洛克（John Locke，1632—1704）清楚地感到货币的充裕会使一切东西变贵，却没有研究这是怎样发生的；而“这项研究的最大困难，在于弄清是通过什么途径”\cite{cantillon1755}。货币增加以后，物价为什么上涨，谁得益，谁受损？新钱进入经济的渠道，为什么还会影响利率？

\subsubsection{渠道}
坎蒂隆的《论一般商业的性质》（\emph{Essai sur la nature du commerce en général}）1755 年才出版，那时他已经去世二十多年。书的第一句话是：“土地是财富的源泉或材料，人的劳动是生产财富的形式。”\cite{cantillon1755}

坎蒂隆先看金银矿。如果新增的货币来自国内的金银矿，“矿主、企业家、熔炼工、精炼工，以及一般所有在那里干活的人，都会按他们的收益增加开支”。他们多买肉、酒、羊毛，物价就上涨。在坎蒂隆看来，多出来的这些开支，必然减少了国内其他居民的份额，这些人起初没有分享到矿山的财富。受损的是谁？“首先是土地的所有者，在租约期内；然后是他们的仆人，以及所有靠固定工资养家的工人”\cite{cantillon1755}。地主的地租在租约到期以前是固定的，工资也是固定的，物价却涨了。新钱从哪里进来，先到谁手里，决定了谁得益、谁受损。这就是后来所说的\textbf{坎蒂隆效应}（Cantillon effect）。

\subsubsection{两个群体的模型}
把坎蒂隆的说法写成最简单的模型。经济里有两个群体：$A$ 先拿到新钱，$B$ 的收入是固定的。实际产出 $Y$ 不变，所有的钱每期都花掉一次（式~\eqref{eq:silver-qty} 里 $V=1$），所以 $P=M/Y$。起初 $A$ 有 $m_A$，$B$ 有 $m_B$，$M=m_A+m_B$。新增的货币 $\Delta M$ 全部到了 $A$ 手里，物价涨到 $P'=(M+\Delta M)/Y$。$B$ 的钱没有变，能买到的产出份额从 $m_B/M$ 降到
\begin{keyeq}
s_B'=\frac{m_B}{M+\Delta M}=\frac{s_B}{1+\Delta M/M},\qquad s_A'=\frac{m_A+\Delta M}{M+\Delta M} .\label{eq:merc-cantillon}
\end{keyeq}
$B$ 的实际购买力按比例 $\Delta M/(M+\Delta M)$ 下降，这部分产出转到了 $A$ 手里。产出没有变，只是重新分了一次。等到 $B$ 的租约与工资重新议定，按新的物价提高收入，份额才会回到原来的样子。所以通胀的分配效应在短期最大，谁的收入调整得慢，谁就吃亏。

\begin{example}{一个月涨 25\%}\label{ex:merc-cantillon}
Velde 说，1720 年 1 月，巴黎的物价“一个月里上涨了 25\%”\cite{velde2003}。（1）一个收入固定的人，这个月的购买力下降了多少？（2）教学假设：$M=100$，其中先拿到新钱的群体 $A$ 有 20，新增的货币 $\Delta M=25$ 全部到了 $A$ 手里，产出不变。按式~\eqref{eq:merc-cantillon}，两个群体的份额怎样变化？（3）从 1720 年 1 月到 8 月，货币增加到 2.6 倍，物价只涨了 19\%\cite{velde2003}。式~\eqref{eq:silver-qty} 里哪些量变了？
\solution （1）$1-1/1.25=20\%$。（2）物价涨 25\%；$B$ 的份额从 $80/100=0.8$ 降到 $80/125=0.64$，实际购买力减少 20\%；$A$ 的份额从 0.2 升到 $45/125=0.36$，实际购买力增加 80\%。（3）由 $MV=PY$ 得 $V/Y=P/M$。$M$ 乘以 2.6，$P$ 乘以 1.19，$V/Y$ 就下降到原来的 $1.19/2.6\approx46\%$：要么流通速度下降了，要么产出增加了，要么两者都有。一种解释是，新增的纸币“绝大部分是大面额的”，多半用来买卖股份或者存着，没有进入日常的买卖，速度下降了；另一种解释是物价指数本身。Velde 提醒，Faure 从法国各地收集了大量证据，虽然多是零星的记载，却显示 1720 年春天物价在上涨，“特别是在商品市场上”，而 Hamilton 的数据“取自医院的记录”\cite{velde2003}。用哪一种价格来衡量，就看到哪一部分人的损失：这正是坎蒂隆的问题。
\end{example}

\subsubsection{1720 年的利率}
新钱的渠道还决定了利率往哪里走。坎蒂隆说：“如果国内充裕的货币是经由放款人之手进来的，它无疑会增加放款人的数目，从而降低通行的利率；但如果它是经由花钱的人进来的，效果就完全相反。”他举的例子正是 1720 年的英格兰。那时货币与流通都很充裕，原来 5\% 或更低的通行利率却没有降下来，反而一直涨到 50\% 与 60\%，因为“人人都成了南海体系里的企业家，都想借钱买股票”\cite{cantillon1755}。钱多了，想借钱的人更多，利率就会上升。今天的货币政策通过银行与债券市场放出货币，首先降低利率；它的分配效应也要看钱先到了谁手里（\chref{ch:heterogeneity}）。

\section{魁奈的经济表}\label{sec:merc-tableau}

坎蒂隆说，土地是财富的源泉；魁奈则认为，只有土地生产剩余。重农学派说的“净产品”是什么？一年的产品与货币怎样在社会的几个阶级之间流转，才能年年重复下去？

\subsubsection{三个阶级}
魁奈画出第一张《经济表》（\emph{Tableau économique}）的时间，有两种说法。杜邦·德·内穆尔（Pierre Samuel du Pont de Nemours，1739—1817）1769 年写道，魁奈“多次对我们说，这一版是 1758 年 12 月的，他有确凿的日期为证”。米拉波侯爵（Victor de Riqueti, marquis de Mirabeau，1715—1789）则说它是 1759 年的，而且不在年初\cite{oncken1888}。最早的几版是一张用锯齿形的线把各项数字连起来的表。1766 年 6 月，魁奈在《农业、商业与财政杂志》上发表了《经济表的算术公式分析》，用几个整数把表的内容讲了一遍，下面的数字都出自这篇文章\cite{quesnay1766}。

魁奈假设一个农业国，土地每年再生产价值五十亿里弗尔的产品。他把国民分成三个公民阶级，即生产阶级、土地所有者阶级与不生产阶级\cite{quesnay1766}。生产阶级是种地的租地农场主与农民，土地所有者阶级包括地主、君主与收什一税的教会，不生产阶级是工匠、制造业者与商人。魁奈称后者“不生产”，不是说他们不干活，而是说他们只是把原料加工成别的形状，卖得的钱只够补偿原料与自己的生活。用他的话说，“他们的预付什么也不生产：花掉了，又回到他们手里，始终作为储备留着”\cite{quesnay1766}。只有土地在补偿了一切成本之后，还留下一个剩余，即\textbf{净产品}（produit net），它以地租的形式归地主所有。

\subsubsection{一年的流转}
魁奈的数字是这样的，单位是十亿里弗尔。生产阶级的\textbf{年预付}（种子、牲畜的饲料、劳动者的生活资料）为 2，生产出 5，其中 2 是净产品或收入。地主有上一年收到的地租 2，其中 1 用于向生产阶级购买，另外 1 用于向不生产阶级购买。不生产阶级的预付 1，用来向生产阶级购买原料\cite{quesnay1766}。不生产阶级的销售合计为 2，其中 1 用于本阶级成员的生活，即向生产阶级买食物，另外 1 留作下一年的预付\cite{quesnay1766}。所以，地主与不生产阶级向生产阶级购买的合计为 3。生产阶级收到的这 3，欠地主 2，作为当年的地租，另外 1 用于向不生产阶级购买制品\cite{quesnay1766}。

产品这边，5 的总产出里，地主与不生产阶级买走了 3。生产阶级还剩 2 的产品，加上从不生产阶级买来的 1 的制品，构成了它每年 3 的基金\cite{quesnay1766}。其中 2 是下一年的年预付，1 用来维持\textbf{原预付}，即农具、牲畜、建筑这类一次投入、多年使用的资本。魁奈假设原预付是年预付的五倍，即 10，那么每年 1 的“利息”就相当于 10\% 的利率\cite{quesnay1766}。他的结论是：“按这个假设，年预付再生产了 250\%。”\cite{quesnay1766}

把每个阶级付出的钱与收到的钱列成一张表，就是表~\ref{tab:merc-tableau}，每一行是谁付钱，每一列是谁收钱。

\begin{table}[htbp]\centering\small
\caption{魁奈的经济表：一年里各阶级之间的货币支付（十亿里弗尔）}\label{tab:merc-tableau}
\begin{tabularx}{\linewidth}{@{}l Y Y Y Y@{}}
\toprule
付款者 \textbackslash{} 收款者 & 生产阶级 & 土地所有者 & 不生产阶级 & 付出合计 \\\midrule
生产阶级 & — & 2（地租） & 1（制品） & 3 \\
土地所有者 & 1（食物） & — & 1（制品） & 2 \\
不生产阶级 & 2（原料 1，食物 1） & — & — & 2 \\\midrule
收到合计 & 3 & 2 & 2 & 7 \\
\bottomrule
\end{tabularx}
\datasource{魁奈（1766）《经济表的算术公式分析》，《农业、商业与财政杂志》第 5 卷第 3 部分，第 11—41 页；计算见 \texttt{checks/check\_ch16.py}。}
\end{table}

设 $T_{ij}$ 为阶级 $i$ 付给阶级 $j$ 的钱，$X$ 为生产阶级的总产出，$A$ 为年预付，$I$ 为原预付的利息，$R$ 为净产品。经济表能年年重复下去，要满足两组条件：
\begin{keyeq}
\sum_{j}T_{ij}=\sum_{j}T_{ji}\quad\text{（每个阶级付出的等于收到的）},\qquad X=A+I+R .\label{eq:merc-tableau}
\end{keyeq}
第一组说，每个阶级一年里收到的钱，正好够它付出去的。表~\ref{tab:merc-tableau} 里每一行的合计都等于同名一列的合计，生产阶级 3，地主 2，不生产阶级 2。第二组说，总产出补偿了年预付与原预付的利息以后，剩下的就是净产品，$5=2+1+2$。这两组条件是\chref{ch:accounts}里部门账户的前身。在部门账户里，一个部门付出的每一笔钱都是另一个部门的收入，各部门的收支差额加起来等于零，全社会的产品与收入相等。在经济表的稳态里，每个阶级的差额恰好都是零，一般的部门账户并不要求这一点。魁奈第一次把一个经济的全年流转画成了一张前后一致的账。

\fig{merc_tableau}{魁奈的经济表（1766 年的算术公式）。看每个方框流进与流出的数是否相等：单位是十亿里弗尔，箭头指向收钱的一方。地主拿上一年的地租 2，向生产阶级买 1 的食物、向不生产阶级买 1 的制品；不生产阶级用预付 1 向生产阶级买原料，卖出制品得到 2，再向生产阶级买 1 的食物，留下 1 作为下一年的预付；生产阶级卖出产品得到 3，付给地主 2 的地租，向不生产阶级买 1 的制品，用来维持原预付。生产阶级 5 的产出中，3 卖给另外两个阶级，2 留作下一年的年预付。示意图，数据：魁奈（1766）。}{fig:merc_tableau}

\begin{example}{经济表的几个数字}\label{ex:merc-tableau}
（1）按表~\ref{tab:merc-tableau}，验证每个阶级付出的等于收到的，以及 250\% 的再生产率与 10\% 的利率。（2）魁奈说，在 20 亿的总收入里，“君主的份额是 5.72 亿；地主的份额是七分之四，即 11.44 亿；什一税征收者的份额是七分之一，即 2.86 亿，都含税”\cite{quesnay1766}。这几个数是怎样算出来的？（3）魁奈区分两种奢侈：“花 100 里弗尔买一升青豆的人，把钱付给了一个种地的人，他把它用在耕种的开支上，有利于每年的再生产；花 100 里弗尔买一条金线饰带的人，把钱付给了一个工匠，工匠要拿出一部分去向外国人买原料。”\cite{quesnay1766}教学假设：地主把 2 亿里弗尔从食物转向用进口原料做的饰带，不生产阶级把多收的 2 亿全部用于向外国购买原料。生产阶级收到的钱变成多少？
\solution （1）生产阶级：收到 $1+2=3$，付出 $2+1=3$；地主：收到 2，付出 $1+1=2$；不生产阶级：收到 $1+1=2$，付出 2。再生产率 $5/2=250\%$；利率 $1/10=10\%$。（2）君主 $2/7$，地主 $4/7$，什一税 $1/7$：$2000\times2/7\approx571.4$，$2000\times4/7\approx1142.9$，$2000/7\approx285.7$（百万里弗尔）。魁奈写作 572、1,144 与 286，三者相加是 2,002，比 20 亿多出 200 万，是他取整的结果。（3）地主向生产阶级买的食物从 10 亿降到 8 亿，不生产阶级多收的 2 亿流到了国外，生产阶级收到的钱从 30 亿降到 28 亿。按经济表的账，它付了 20 亿地租以后，只剩 8 亿去买制品维持原预付，或者要从年预付里挪钱；不管哪一种，下一年的再生产都会缩小。魁奈由此主张，地主的花费要多流向土地的产品。
\end{example}

\subsubsection{政策}
重农学派的政策从这张表里推出来。既然只有土地产生净产品，一切税最终都落在净产品上，那么“地主、君主与整个国家都有极大的利益，让赋税全部直接征在土地的收入上”\cite{quesnay1766}。经济表的数字还假定，自由贸易维持着产品的好价格，例如每塞蒂埃小麦 18 里弗尔\cite{quesnay1766}。塞蒂埃（septier）是法国量粮食的容量单位。粮价要好，种地才有利可图，所以粮食的买卖应当自由。重商主义者盯着金银的进出，重农学派盯着土地的收入，两者开出的药方于是很不一样。

\begin{compare}{重商主义者与重农学派}{重商主义者（孟）}{重农学派（魁奈）}
\cmprow{财富是什么}{金银，贸易顺差带回来的“财宝”}{土地每年生产的净产品}
\cmprow{谁是生产的}{对外贸易，以及为出口加工的制造业}{只有农业；制造业与商业是“不生产的”}
\cmprow{货币}{金银就是“财宝”；但要用在贸易里，孟也承认钱多了会抬高物价、减少贸易}{流通的媒介，自己不生产任何东西}
\cmprow{政策}{多卖少买，限制进口，鼓励出口}{粮食自由贸易、维持粮价；赋税只征在土地的收入上}
\cmprow{分析的单位}{一国与外国之间的收支}{一国之内三个阶级之间的循环流转}
\cmpsum{两派都在找一国财富的源头：重商主义者把它放在国界上，重农学派把它放在土地上；魁奈留下的不是“只有农业生产”的结论，而是把全年的流转画成一张前后一致的账的方法。}
\end{compare}

\section{数据检验与反例}\label{sec:merc-test}

斯密说，密西西比计划也许是世界上从未见过的最荒唐的计划。后来的经济史学家要问的是另一件事。劳的公司与南海公司的股价，是不是离开了任何合理的估值？买股票的人是不是都失去了理智？

\subsubsection{印度公司值多少}
1719 年 12 月宣布的每股 200 里弗尔股息，公司挣得出来吗？Velde 把当时人与自己对公司每年收益的估计排在一起（表~\ref{tab:merc-earnings}）。劳的两个估计是 9,100 万与 8,000 万里弗尔，Dutot 的是 8,050 万，Velde 修订后为 7,550 万\cite{velde2003}。按 Velde 的数字，把收益平分给 60 万股，每股股息只有约 125 里弗尔，$7550\text{ 万}/600{,}000\approx125.8$。

\begin{table}[htbp]\centering\small
\caption{印度公司预期的年收益（百万里弗尔，1719 年 12 月）}\label{tab:merc-earnings}
\begin{tabularx}{\linewidth}{@{}l Y Y Y Y@{}}
\toprule
来源 & 劳 (1) & 劳 (2) & Dutot & Velde 修订 \\\midrule
国王的债务（3\% 的利息） & 48 & 48 & 48 & 48 \\
总包税所 & 12 & 8 & 15 & 10 \\
直接税 & 1 & 1 & 1.5 & 1 \\
铸币 & 12 & 10 & 4 & 0 \\
烟草 & 6 & 5 & 2 & 10 \\
贸易 & 12 & 8 & 10 & 6.5 \\\midrule
合计 & 91 & 80 & 80.5 & 75.5 \\
\bottomrule
\end{tabularx}
\datasource{Velde（2003）表 10。劳 (1)、劳 (2) 是劳本人的两种估计。}
\end{table}

收益值多少钱，要看投资人要求多高的收益率，以及收益会不会增长。\chref{ch:finrevolution}的永续年金公式，式~\eqref{eq:finrev-perp}，假定每年的支付不变。如果股息每年按 $g$ 增长，从 $d$ 开始，按收益率 $r$ 贴现（$r>g$），股价就是
\begin{keyeq}
P=\frac{d}{r-g} .\label{eq:merc-gordon}
\end{keyeq}
这叫戈登公式（Gordon formula）。$g=0$ 时，它就是式~\eqref{eq:finrev-perp}。它说明，同样的股息，$r-g$ 越小，估值越高；$r-g$ 从 3\% 降到 1.5\%，股价就翻一番。

\begin{example}{三种估值}\label{ex:merc-gordon}
（1）Velde 说，1723—1769 年，印度公司的价格股息比在 4 到 24 之间波动，“平均约为 15”。以这个倍数乘以 7,550 万里弗尔的收益，公司值多少？每股值多少？与 9,000 里弗尔相比如何？（2）公司 1719 年底的市值约为 49.6 亿里弗尔。按式~\eqref{eq:merc-gordon}，7,550 万的收益要多大的 $r-g$ 才值这么多？只假定利率降到荷兰的 3\%、没有增长呢？（3）每股 125 里弗尔的股息，按 3\% 的利率值多少？如果再加上 0.5\% 的增长呢？
\solution （1）$75.5\times15\approx1{,}132$ 百万里弗尔；Velde 说，这相当于每股 1,875 里弗尔，“是 9,000 里弗尔顶点股价的五分之一”，也就是说股价是这个估值的五倍\cite{velde2003}。$1132\text{ 百万}/1875\approx60$ 万股，与股数相符。（2）$r-g=75.5/4957.5\approx1.5\%$。Velde 说，要证明市值合理，需要“荷兰那样 3\% 的利率与 1.5\% 的增长率，在工业革命以前，没有一个欧洲国家有过这样的增长”；只假定利率降到荷兰的水平，$75.5/0.03\approx2{,}517$ 百万，市值仍是这个估值的约两倍（$4957.5/2517\approx1.97$）\cite{velde2003}。作为比较，1726—1789 年法国的财政收入实际每年增长约 0.6\%，GDP 增长约 0.5\%\cite{velde2003}。（3）$125/0.03\approx4{,}167$ 里弗尔；加上 0.5\% 的增长，$125/(0.03-0.005)=5{,}000$ 里弗尔，正是劳 1719 年 9 月发行新股、启动债务置换的价格\cite{velde2003}。
\end{example}

Velde 写道：“很难避免这样的结论：公司被高估了好几倍。”可他也说：“凭着现代的理论与经验，再加上相当的乐观，可以承认这个体系本来是可能成功的。”\cite{velde2003}

关键在于价格怎样形成。在 Velde 看来，股价的上涨本身是先暗中、后公开托市的结果，部分是为了引诱债主自愿接受债务置换。他说，“最晚从 1720 年 1 月起，也许从 1719 年 11 月或 12 月起，所谓的‘市场’价格只能看作劳的政策”\cite{velde2003}。图~\ref{fig:merc_shares}(a) 的褐色虚线说明了这一点。钉住期间名义股价不变，按白银计的股价从 3 月就开始下跌；到 1720 年 11 月，它回到了 1718 年的水平。Velde 说，1722 年公司改组以后，考虑到股份的拆分与记账单位的变化，股价与 1718 年大致相同\cite{velde2003}。

\subsubsection{南海公司值多少}
Garber 对南海公司做了同样的计算。1720 年 9 月初，南海公司股份的市值为 1.64 亿英镑。支撑这个价格的看得见的资产，是公司对政府的债权带来的收入，1727 年以前每年 190 万英镑，此后每年 150 万英镑。按 4\% 的长期贴现率，这笔资产约值 4,000 万英镑\cite{garber1990}。按年金与永续年金的公式，$1.9\times[1-1.04^{-7}]/0.04+(1.5/0.04)/1.04^{7}\approx11.4+28.5=39.9$ 百万英镑。扣掉公司答应付给政府的 710 万与欠下的 600 万债券和票据，净资产为 2,610 万英镑。另外还有借给股东的 1,100 万，以及现金认购者将来要付的 7,000 万应收款。Garber 说，股份的价值超过资产 6,000 多万英镑；考虑到现金债权靠不住，股份的价值是有形净资产的五倍或更多\cite{garber1990}。

Garber 并不因此认为买股票的人疯了。他问，什么样的无形资产能证明公司值这么多？他的回答仍是劳的预言，即随着一笔“信用基金”的积累，商业会扩张。他还引了 Scott 的看法，400 英镑的价格并不过分。回头看，凡是预期商业回报高到足以支撑更高股价的人，大概都太乐观了。可是在 Garber 看来，“这段插曲完全可以理解为投机者根据当时最好的经济分析行事，随着他们对市场基本面的看法改变而推动价格”\cite{garber1990}。他批评经济学家迷恋于宏大投机事件的光彩，挤在泡沫的解释里，忽视了对可能的基本面的考察\cite{garber1990}。

\subsubsection{横截面的证据}
如果 1720 年的投资人只是盲目追涨，所有的股票应当一起涨。Frehen、Goetzmann 与 Rouwenhorst 收集了伦敦、巴黎与荷兰各城市的每日股价，发现情形不是这样。他们找到的证据，反驳了“不加区别的非理性亢奋”，却支持围绕两个因素的投机，一是大西洋贸易，二是保险公司的设立\cite{fgr2009}。北欧四家做大西洋贸易的公司，即密西西比公司、皇家非洲公司、南海公司与荷兰西印度公司，在这场全球泡沫里都涨了 7 到 10 倍；两家做亚洲贸易的东印度公司涨得少得多\cite{fgr2009}。图~\ref{fig:merc_shares}(b) 里南海公司与东印度公司的差别就是一例。

这里的“大西洋贸易”离不开奴隶贸易。南海公司持有向西属美洲贩运非洲奴隶的专营权，皇家非洲公司同样做大西洋的奴隶贸易。三角贸易把工业品运到非洲，把非洲人作为奴隶运到新大陆，再把种植园的产品运回欧洲，作者称它为“近代早期最重要的国际经济制度之一”。这种贸易在 1720 年才刚刚开始，利润还远在将来，可在作者看来，也许正是这种未来的前景，激起了当时投资人的热情\cite{fgr2009}。投资人押注的这门生意，建立在被贩卖的非洲人的自由与生命之上。

他们的表 1 依据 1721 年的一幅讽刺版画《泡沫者之镜》（The Bubbler's Mirror）。版画列出了伦敦各公司股份的面值，以及泡沫期间超出面值的最大百分比\cite{fgr2009}。按这张表，伦敦保险公司相对面值的最大涨幅平均为 2,013\%，大西洋贸易公司为 895\%。太平洋（东印度公司）是 349\%，银行与金融是 335\%，全部 45 家平均 1,172\%。荷兰的公司平均只有 127\%\cite{fgr2009}。

荷兰提供了一个对照。荷兰没有一家公司模仿密西西比公司或南海公司的金融操作，荷兰的股价却也涨了，包括西印度公司在内的其他类型的公司都出现了泡沫。作者由此推断，政府的债务转换“不可能是英国泡沫的唯一基础”，密西西比与南海泡沫主要不是由于对债务—股份置换这种金融创新的投机\cite{fgr2009}。伦敦的两家保险公司在同一天达到顶点，即 8 月 26 日（格里历），第二天开始下跌。南海公司的崩溃从 9 月 19 日（格里历）开始。他们认为，引发崩溃的首要原因看来是政府恰在此时援引《泡沫法》、把英国公司限制在特许状范围内的令状。崩溃一来，不管投资人的预期建立在大西洋贸易、海上保险还是房地产交易上，9 月与 10 月所有的股票都很惨\cite{fgr2009}。两家保险公司后来一直存续到现代，证明了股份制保险公司的长期生命力\cite{fgr2009}。

\subsubsection{明知是泡沫，也要乘上去}
Temin 与 Voth 研究了霍尔银行（Hoare's Bank）在 1720 年的投资。这是霍尔家族的私人银行，创办人原是金匠，1690 年迁到舰队街以后才专做银行业务\cite{tv2004}。他们认为，这家银行“知道泡沫正在形成，却仍然投资于这只股票；‘乘着泡沫’是有利可图的”\cite{tv2004}。霍尔银行赚了 28,000 多英镑。当时一个中产家庭每年有 200 英镑就过得很舒服了\cite{tv2004}，这笔利润相当于这样一个家庭 140 年的收入。在他们看来，可以预见的投资者情绪能阻止人们攻击泡沫；理性的投资者也许只在某个协调事件促成共同行动时，才会发动攻击\cite{tv2004}。

也有研究认为，1720 年的投资人确实不理性。南海公司的股票与它的现金认购凭证（subscription receipts）同时在市场上交易，两者可以互相替代。Dale、Johnson 与 Tang 用新找到的这两个平行市场的价格，发现两者之间的估值差距越来越大，理性的解释并不充分。他们把这些反常归因于投资人有偏差的决策，以及他们“把金融市场看作赌博市场的倾向”\cite{djt2005}。Shea 反驳说，认购合同里含有违约的选择权，这在 18 世纪的认购融资里很常见。他用公司的记录与当时的小册子说明，人们懂得法律规定的这套融资办法，并认为 Dale、Johnson 与 Tang 的分析“有无法挽回的缺陷”\cite{shea2007}。

劳的认购凭证也是这样。误了一期分期款，就失去股份与以前付的全部款项，所以 Velde 说，这些凭证是股份的期权，而不是股份\cite{velde2003}。争论的关键是两种证券能否互相替代。如果认购凭证里含有一份期权，它的价格本来就该与股票不同，比较两者以前，要先给这份期权定价。

\begin{pitfall}[泡沫说明买股票的人都失去了理智]
本节的大部分证据并不支持这个说法。Velde 的计算说明，按任何合理的利率与增长率，9,000 里弗尔的印度公司股份都被高估了；但价格是劳托着的，买股票的人看到的是公司按 9,000 里弗尔回购的承诺。Garber 说南海公司的投资人是在按当时最好的分析下注。Frehen、Goetzmann 与 Rouwenhorst 发现涨得最多的是做大西洋贸易与保险的公司，不是所有股票一起涨。Temin 与 Voth 发现，明知股价过高的人也可能理性地买进，因为他们预期能在别人之前卖掉。Dale、Johnson 与 Tang 认为投资人确有偏差，但他们的证据本身还在争论之中。“泡沫”描述的是价格离开了基本面，并不说明每个人都不理性；要问的是，在那样的制度与承诺之下，什么样的预期让价格走到了那里。
\end{pitfall}

\begin{deeper}[史料与几处不一致]
\textbf{南海公司的价格与日期。}顶点的 950 英镑来自 Neal，Scott 记为 1,050 英镑，Garber 说后者大概包括了 10\% 的股票股息；本章的数据（Castaing 的行情表，经 Frehen、Goetzmann 与 Rouwenhorst 整理）在儒略历 6 月 29 日至 7 月 1 日为 950 英镑。Garber 写 8 月 31 日约 775 英镑，同一数据在这一天为 810 英镑。Garber 说 6 月 17 日的认购以后，股价“立刻从 745 英镑跳到 950 英镑”；在这份数据里，6 月 17 日的价格是 755 英镑，到 22 日也只有 765 英镑（745 英镑是 6 月 10 日的价格）；23 日到 28 日没有报价，950 英镑最早出现在 6 月 29 日，那时过户登记已经关闭。公司答应付给政府的钱，Garber 在一处说是“至多 750 万英镑”，条件是收进全部 3,100 万国债；算净资产与讲善后时，用的却是 710 万。伦敦当时用儒略历，巴黎与荷兰用格里历，两者相差 11 天；正文里的伦敦日期按原文（儒略历），图按格里历。

\textbf{1716 年的银行。}Velde 说通用银行设立于 1716 年 5 月，Garber 写的是 6 月。Velde 在第 9 页说国家票据 1716 年年中折价 37\%，在第 11 页说银行发行股份时折价 60\%；“只要 690 里弗尔现金”与 60\% 的折价相符，按 37\% 算则约为 900 里弗尔。

\textbf{Velde 的几处笔误。}Velde 说体系以前国家是每年约 9,000 万里弗尔的“债权人”（creditor），从上下文看是债务人；又说 5 月 21 日的法令要把“股份”降到面值的一半，第二个“股份”应当是纸币；正文里把 7 月发行的股份叫作“女儿”，可脚注 19 与同一句的下文都表明“女儿”是 6 月的第二次发行，7 月的是“孙女”，本章按脚注。

\textbf{劳的几个数字。}Garber 说钉住期间的四次增发为 3 亿、3.9 亿、4.38 亿与 3.62 亿里弗尔，合计 14.9 亿，“大约一个月里”流通的纸币翻了一番；Velde 说公司 3—5 月花了 12.135 亿收购股份，表 4 的纸币同期增加了 11.452 亿。1720 年 1 月的物价，Velde 说上涨了 25\%，Garber 说是 23\%，两人用的都是 Hamilton 的指数。Visa 把 22.112 亿核减到 17 亿，按这两个数算是 23.1\%，Velde 写作 23.5\%。Velde 正文说王家银行的纸币“到 1719 年 7 月”为 4 亿、“1720 年 1 月”达到 10 亿，表 4 却是 1.599 亿（1719 年 7 月 31 日）与 8.101 亿（1720 年 1 月 30 日），到 8 月底与 2 月底才接近这两个数；本章正文与图~\ref{fig:merc_notes} 用表 4。

\textbf{估值。}Velde 说 125 里弗尔的股息与 3\% 的利率能让股价“达到约 5,000 里弗尔”；按式~\eqref{eq:merc-gordon} 不计增长只有 4,167 里弗尔，要再加 0.5\% 的增长才是 5,000。Garber 的“约 4,000 万英镑”按年末付款、从 1720 年起贴现约为 3,990 万；减去 710 万与 600 万应为 2,680 万，原文写作 2,610 万；按原文列出的数，$164-26.1-11-70\approx56.9$ 百万英镑，略低于原文说的“6,000 多万”，差别大概来自贴现的时点与所用的股价。

\textbf{魁奈的表。}第一张经济表的日期，魁奈说是 1758 年 12 月，米拉波说是 1759 年；本章用 1766 年的《算术公式分析》，它比最早的表简单。魁奈的三份税收相加多出 200 万里弗尔（例~\ref{ex:merc-tableau}）。

\textbf{作者与版本。}孟的书约 1630 年写成，1664 年才出版，斯密把书名引作 \emph{England's Treasure in Foreign Trade}。坎蒂隆的书 1755 年出版时他已去世，本章用的维基文库本把作者标为“Richard Cantillon (attribué)”，即“认为是”坎蒂隆所作。印度公司的股价，1719 年 8 月以前来自《阿姆斯特丹报》，此后来自 Murphy 的整理，钉住期间固定的报价换成了莱顿报纸上的市场价格（据数据库的说明）；1720 年的价格以不断贬值的纸币计，所以图~\ref{fig:merc_shares}(a) 另画了按白银折算的价格。
\end{deeper}

\section{今天还剩什么}\label{sec:merc-legacy}

劳的银行与南海公司早已不在了。孟与魁奈的一些想法后来被推翻了。孟把金银当作财富，可金银不等于财富（\secref{sec:merc-mun}）。魁奈把工匠、制造业者与商人归入不生产阶级，斯密说，重农学派体系的根本错误，似乎就在于把这些人说成“完全不生产的”\cite{smith1776}（\chref{ch:smith}）。另外几个问题今天还在。

\textbf{贸易顺差的执念。}“卖多买少”仍然是常见的政策目标，对某一国的双边逆差仍然常被说成损失。孟的四个漏洞里，世界的顺差之和为零、顺差会自我纠正这两个，今天换成了全球失衡与贸易冲突的问题（\chref{ch:fragmentation}）。顺差怎样通过物价与汇率自我纠正，是\chref{ch:bullion}与\chref{ch:external}的内容。

\textbf{中央银行买国债。}坎蒂隆认为可以确定，一家与大臣有默契的银行，能够按这位大臣的意愿抬高并支撑公债的价格，压低国内的利率。可是，“如果某种恐惧或意外促使持票人向银行要钱，炸弹就会爆开，人们会看到这是危险的操作”\cite{cantillon1755}。今天的央行大量购买国债以压低长期利率（\chref{ch:zlb}）；货币政策是不是在替财政融资，是\chref{ch:fiscaldominance}的问题。劳的体系是一个极端的例子。1720 年初，发钞银行并入了包税的公司，公司的首脑又当上了财政大臣。

\textbf{钉住与货币数量。}按固定价格无限量地买卖某种资产，就放弃了对货币数量的控制，式~\eqref{eq:merc-peg}。钉住汇率与资本流动的三难选择，是\chref{ch:mundell}的问题。

\textbf{分配效应。}新钱先到谁手里，决定了通胀与资产价格上涨的分配后果，式~\eqref{eq:merc-cantillon}。今天的异质主体模型研究货币政策怎样通过不同人的收入、资产与负债起作用（\chref{ch:heterogeneity}）。

\textbf{政府股权。}体系里出现的另一个概念是政府股权，即政府用一种收益随经济状况变化的凭证，代替固定利息的国债。Velde 说，这个概念还没有被正式地重新采用\cite{velde2003}。债务能否维持，是\chref{ch:budget}的问题。

\textbf{劳的名声。}Velde 说，到 20 世纪 30 年代，随着人们对不兑现纸币有了更多经验，并且相信政府政策（包括货币政策）可以而且应当用来刺激经济，“劳的名声上升了”\cite{velde2003}。有失业时多发货币能否增加产出，是\chref{ch:generaltheory}的问题。

\textbf{循环流量与投入产出。}魁奈的表是国民核算（\chref{ch:accounts}）的前身。在两者的账上，一个部门的支出都是另一个部门的收入，产出补偿了中间投入与资本的维持以后，剩下的就是收入。马克思的再生产图式与后来的投入产出表都从这张表出发（\chref{ch:marx}）。

\begin{takeaways}
\begin{itemize}
\item 孟（约 1630 年写成，1664 年出版）认为财富是贸易顺差带回来的金银，$\Delta G=X-M$：出口 220 万、进口 200 万英镑，王国每年富 20 万英镑。他区分了王国的收益与商人的利润，为东印度公司出口白银辩护；但他把金银当作财富，把国内的资源当作免费的，而且世界的顺差之和为零，顺差还会自我纠正（\chref{ch:bullion}）。
\item 劳（1705）认为货币的多少决定能雇多少人，信用就是货币，纸币可以以土地为担保、与需求相等。在有闲置资源时，这是 $MV=PY$ 里 $Y$ 随 $M$ 增加的情形；危险在于“需求”可以无限增长。
\item 劳的体系把平均利率 4.5\% 的国债换成公司的股份。按式~\eqref{eq:merc-swap}，市盈率高于 $1/0.045\approx22$ 时，置换对国王合算；1719 年 12 月，价格股息比为 47。1720 年 3 月 5 日起，公司按 9,000 里弗尔钉住股价，$\Delta M=\bar p\,\Delta S$，式~\eqref{eq:merc-peg}，3—5 月为收购股份花了 12.135 亿里弗尔。纸币从 3,800 万增加到 23.59 亿里弗尔；5 月 21 日的法令以后，纸币按白银计的价值在两个月里跌掉三分之二，纸币持有人最后平均拿回面值的约 20\%。
\item 南海公司取得了收购约 3,100 万英镑国债的权利，收进了其中的大部分；给长期年金持有人的报价按 400 英镑的股价约值 3,375 英镑，年金本身约值 1,600 英镑。股价 7 月达到 950 英镑，8 月 31 日到 10 月 1 日从约 775 跌到约 290 英镑；9 月初 1.64 亿英镑的市值，对应约 4,000 万英镑的看得见的资产。Velde 估计，按历史的价格股息比，印度公司 9,000 里弗尔的股价约为其价值的五倍，即使利率降到荷兰的水平也约为两倍；Frehen、Goetzmann 与 Rouwenhorst 发现涨得最多的是大西洋贸易与保险公司，Temin 与 Voth 发现有人明知是泡沫也乘了上去。
\item 坎蒂隆追问新钱先到谁手里：先拿到的人得益，收入固定的人按 $\Delta M/(M+\Delta M)$ 的比例受损，式~\eqref{eq:merc-cantillon}；钱经由借款的人进来时利率会上升，1720 年伦敦的利率涨到 50\%—60\%。魁奈的经济表把年预付 20 亿、产出 50 亿（再生产 250\%）、净产品 20 亿的流转画成一张账，每个阶级付出的等于收到的，式~\eqref{eq:merc-tableau}，是国民核算与投入产出表的前身。
\end{itemize}
\end{takeaways}

\begin{selfcheck}
\begin{enumerate}
\item 一国每年出口货物 300 万英镑，进口货物 340 万英镑，本国的船为外国人运货，每年收运费 30 万英镑。只算货物，金银流入多少？像孟编全国的贸易差额那样把运费算进去呢？如果它的贸易伙伴也都照孟的规则做，会发生什么？
\item 国王欠债 10 亿里弗尔，利率 5\%。一家公司发行 20 万股换掉这笔债，公司每年的收益为 4,000 万里弗尔。按式~\eqref{eq:merc-swap}，股价至少要多高，置换对国王才合算？如果股价是 5,000 里弗尔，国王的借款成本是多少？
\item 一家公司按每股 8,000 里弗尔钉住股价，原来流通的纸币为 10 亿里弗尔。公众一个月里卖给公司 5 万股。按式~\eqref{eq:merc-peg}，纸币增加多少？如果货币的流通速度与产出都不变，物价会涨多少？
\item 一笔年金的市价为 1,500 英镑，公司的报价是 6 股加 500 英镑现金。股价跌到多少时，接受报价会吃亏？股价为 300 英镑时，报价值多少？
\item 每股股息 150 里弗尔，投资人要求 4\% 的收益率，股息每年增长 1\%。按式~\eqref{eq:merc-gordon}，股价是多少？如果股息不增长呢？如果市价是 9,000 里弗尔，隐含的 $r-g$ 是多少？
\item 按式~\eqref{eq:merc-cantillon}，设 $M=200$，先拿到新钱的群体有 50，新增的 50 全部到了这个群体手里。收入固定的群体的实际购买力下降多少？先拿到新钱的群体的份额变成多少？
\end{enumerate}
\end{selfcheck}

\begin{answers}
\item 只算货物，$\Delta G=300-340=-40$ 万英镑，金银流出。算上运费，$-40+30=-10$ 万英镑，仍然流出，但少得多。孟的账已经把运费、保险费与商人的利润算进了出口；今天的国际收支表把这些收入记在服务贸易里（\chref{ch:external}）。所有国家的贸易差额之和为零，不可能人人顺差；大家都限制进口、鼓励出口，结果是贸易缩小，而不是人人得到金银。
\item 每股收益 $E/N=4000\text{ 万}/200{,}000=200$ 里弗尔；要 $P/(E/N)>1/0.05=20$，股价要高于 4,000 里弗尔。按 5,000 里弗尔，20 万股值 10 亿，正好换掉全部债务，成本为 $4000\text{ 万}/10\text{ 亿}=4\%$，低于 5\%。
\item $\Delta M=8000\times50{,}000=4$ 亿里弗尔，纸币增加 40\%。按式~\eqref{eq:silver-qty}，$V$ 与 $Y$ 不变时物价也涨 40\%。
\item $6P+500=1500$，$P\approx166.7$ 英镑；股价为 300 英镑时，报价值 $6\times300+500=2{,}300$ 英镑。
\item $P=150/(0.04-0.01)=5{,}000$ 里弗尔；不增长时 $150/0.04=3{,}750$ 里弗尔。市价 9,000 里弗尔时，$r-g=150/9000\approx1.7\%$。
\item $\Delta M/(M+\Delta M)=50/250=20\%$；收入固定的群体的份额从 $150/200=0.75$ 降到 $150/250=0.6$。先拿到新钱的群体的份额从 $50/200=0.25$ 升到 $100/250=0.4$。
\end{answers}

回到劳当上财政大臣的那一天。他管着法国唯一的发钞银行，他的公司替国王收税，也替国王还债。他相信钱不必是金银，钱多了，人就有活干。不到一年，他离开了法国，纸币持有人后来平均只拿回面值的约五分之一。劳在巴黎印钞的时候，欧亚大陆的另一端也在为货币发愁。德川日本按米计税，却要管理金、银、铜三种货币；清代中国的白银与铜钱，比价不停地变。一个以米计税的国家，怎样管理三种钱与物价？
